\chapter{The Constraint Vocabulary}\label{chap:7}

\begin{quote}
\emph{``The limits of my language mean the limits of my world.''} ---
Ludwig Wittgenstein
\end{quote}

Constraints are the verbs of a CP-SAT model. Where \Cref{chap:6} showed how
to declare the unknowns, this chapter shows how to bind them --- every
relation, restriction, and rule that distinguishes a valid configuration
from an invalid one. CP-SAT ships with a curated set of constraint
families. Some \emph{global} families in the catalogue carry a dedicated
propagator far stronger than what a modeler would write out of plain
inequalities; others are flattened internally to clauses or linear
constraints. Either way, calling the global by name expresses intent and
lets the solver pick the best encoding.

\section{\texorpdfstring{Linear Constraints --- \texttt{model.add}}{Linear Constraints --- model.add}}\label{sec:7.1}

The bread and butter of any CP-SAT model.
\lstinline!model.add(...)! accepts any linear (in)equality
and registers it as a constraint.

\begin{lstlisting}[language=Python]
model.add(10*x + 15*y <= 200)
model.add(x + z == 2*y)
model.add(x < y + z)  # rewritten internally to x <= y + z - 1
# Double-bounded:
model.add_linear_constraint(10*x + 15*y, lb=-100, ub=10)

# Variable-length sums are built with Python's built-in sum(...):
model.add(sum(w[i] * take[i] for i in items) <= capacity)
# Exactly-one over a row of Booleans:
model.add(sum(x[i, j] for i in F) == 1)
\end{lstlisting}

The expressions on the left are built by overloading:
\lstinline!+!, \lstinline!*!,
\lstinline!-!, \lstinline!==!,
\lstinline!<=!, \lstinline!<!,
\lstinline!>=!, \lstinline!>! on CP-SAT
variables (and integer constants) all return CP-SAT expression or
constraint objects. Python's built-in \lstinline!sum(...)!
chains the additions for you, so a generator expression over a family of
variables --- the workhorse pattern for capacity and counting
constraints --- drops directly into
\lstinline!model.add(...)!. The constraint object that
comes out is registered with the model.

CP-SAT works exclusively in integers, so strict inequalities are
rewritten to non-strict-with-shift: \(x < y\) becomes \(x \le y - 1\).
The semantics are identical, but knowing that the rewrite happens
prevents confusion when reading the solver's logs or inspecting an
exported model. For double-bounded expressions --- \emph{the load is
between 50 and 100} ---
\lstinline!add_linear_constraint(expr, lb, ub)! is
occasionally cleaner than two separate
\lstinline!model.add! calls; both forms produce the same
model.

The shapes that recur in essentially every model are capacity
(\(\sum_i w_i x_i \le C\)), conservation laws (\emph{inflow equals
outflow plus storage}), counting (\emph{exactly k of these are taken}),
budget constraints, and flow conservation per node
(\(\sum_{j \in \delta^+(i)} x_{ij} - \sum_{j \in \delta^-(i)} x_{ji} = b_i\)).

With practice, almost any constraint can be \emph{expressed or
approximated} as a linear (in)equality --- often with a few helper
Booleans or auxiliary variables. Mixed-Integer Programming solvers
(Gurobi, CPLEX, HiGHS) accept \emph{only} this family of constraints,
and entire industries run on models built from
\lstinline!model.add(...)! alone. CP-SAT exposes a much
richer vocabulary in the sections ahead --- Boolean logic, reification,
scheduling globals, routing globals --- but the linear shape is the
universal one. The richer constraints are \emph{convenience and
propagation strength}, not necessity. Understanding linear-only modeling
is therefore a transferable skill: it carries over to every MIP solver
and remains the lingua franca even when CP-SAT's globals are the better
choice for a given problem.

\section{\texorpdfstring{Boolean Logic --- \texttt{add\_bool\_or}, \texttt{add\_bool\_and}, \texttt{add\_implication}}{Boolean Logic --- add\_bool\_or, add\_bool\_and, add\_implication}}\label{sec:7.2}

Pure logical constraints over Boolean \emph{literals} --- Booleans or
their negations.

\begin{lstlisting}[language=Python]
model.add_bool_or([a, b, c])       # at least one true
model.add_at_least_one([a, b, c])  # alias
model.add_at_most_one([a, b, c])
model.add_exactly_one([a, b, c])
model.add_bool_and([a, ~b, c])     # all true (mostly inside only_enforce_if)
model.add_implication(a, b)        # a => b  (binary only)
model.add_bool_xor([a, b, c])      # odd number true
\end{lstlisting}

The arguments are \emph{literals}: Boolean variables or their negations
via \lstinline!\~!. An integer variable with domain
\(\{0, 1\}\) is \emph{not} a literal and will be rejected --- if the
model needs logical operations, declare the variable with
\lstinline!new_bool_var! from the start.

CP-SAT does not let you nest Boolean constraints --- each call takes
\emph{literals}, never formulas. In modeling, however, the rules you
want to express rarely arrive in flat form. A typical specification
might read \emph{``at least one of three production lines is fully
prepared''} --- formally
\((a_1 \land b_1) \lor (a_2 \land b_2) \lor (a_3 \land b_3)\), a
disjunction of conjunctions --- or \emph{``if either of two combined
conditions holds, ship today''}, mixing conjunction, disjunction, and
implication in a single formula. CP-SAT will not accept these as
written. Unfolding the formula into flat clauses is part of the modeling
work, and it is \emph{classical propositional logic} --- the same
machinery any first course on logic introduces: De Morgan's laws,
distribution, conjunctive normal form, and Tseitin's transformation.

Two routes are available, and which one fits depends on the size and
shape of the formula:

\paragraph{Direct unfolding by distribution.} Push negations inward via De
Morgan (\(\lnot(x \land y) \equiv \lnot x \lor \lnot y\),
\(\lnot(x \lor y) \equiv \lnot x \land \lnot y\)) and distribute
disjunction over conjunction
(\(x \lor (y \land z) \equiv (x \lor y) \land (x \lor z)\)) until the
formula is in \emph{conjunctive normal form} --- a conjunction of
clauses, where each clause is a disjunction of literals. Each clause
maps to one \lstinline!add_bool_or!. Distribution can
blow up the formula size: a disjunction of \(n\) binary conjunctions has
up to \(2^n\) clauses in CNF, which is fine for a few cases but quickly
impractical.

\paragraph{Tseitin's transformation with auxiliary Booleans.} When direct
distribution would be too large, introduce a fresh Boolean \(t_S\) for
every sub-expression \(S\) and constrain it with
\(t_S \Leftrightarrow S\) (three short clauses for a binary conjunction
or disjunction). The compound formula then refers to the auxiliaries in
place of their sub-expressions, and the encoding grows \emph{linearly}
in the size of the formula. This is standard SAT preprocessing, in
widespread use since Tseitin's 1968 paper.

A worked example, on the formula above:

\((a_1 \land b_1) \;\lor\; (a_2 \land b_2) \;\lor\; (a_3 \land b_3)\)

\emph{Direct CNF.} Distribute disjunction over the three conjunctions:
pick one literal from each, giving \(2^3 = 8\) clauses such as
\(a_1 \lor a_2 \lor a_3\), \(a_1 \lor a_2 \lor b_3\), \ldots,
\(b_1 \lor b_2 \lor b_3\). Eight \lstinline!add_bool_or!
calls. With four pairs the count rises to sixteen, with eight pairs to
256 --- direct distribution does not scale.

\emph{Tseitin form.} Introduce \(t_1, t_2, t_3\) with each
\(t_i \Leftrightarrow a_i \land b_i\), then assert the disjunction over
the auxiliaries:

\begin{lstlisting}[language=Python]
t = [model.new_bool_var(f"t_{i}") for i in range(3)]
for i in range(3):
    # t_i <=> (a_i AND b_i): three clauses
    model.add_implication(t[i], a[i])           # t_i => a_i
    model.add_implication(t[i], b[i])           # t_i => b_i
    model.add_bool_or([~a[i], ~b[i], t[i]])     # a_i AND b_i => t_i
model.add_bool_or(t)                            # t_1 OR t_2 OR t_3
\end{lstlisting}

Three auxiliary Booleans, ten clauses total, scaling linearly with the
number of pairs. The two encodings are logically equivalent on the
original variables; they differ in how they exploit the solver's
machinery.

\lstinline!only_enforce_if! (\Cref{sec:7.3}) automates a narrow
slice of this --- specifically, conditional enforcement of a single
constraint under a literal or short conjunction of literals. It does not
generalize to arbitrary nested formulas; for those, the manual unfolding
above is the standard approach.

A concrete consequence for the basic vocabulary above:
\lstinline!add_implication! is binary only
(\(a \Rightarrow b\)). Implications with several antecedents
(\(a \land b \Rightarrow c\)) flatten by definition of implication to
\(\lnot a \lor \lnot b \lor c\) --- a single call to
\lstinline!add_bool_or([\~a, \~b, c])! --- without
needing an auxiliary.

Three Boolean patterns appear in essentially every Boolean-rich model:
exactly-one over an assignment row --- \emph{each customer goes to
exactly one warehouse}, the FLP assignment constraint from \Cref{sec:5.5} ---
at-most-one for mutually exclusive flags such as pricing tiers or
compatibility classes, and binary implication for \emph{if A then B}
rules.

Short clauses --- binary \lstinline!add_bool_or([a, b])!
in particular --- are among the cheapest propagation CP-SAT performs,
and the SAT layer's clause-learning machinery tends to do well when the
encoded logic decomposes into many small clauses rather than a few wide
ones.

\section{\texorpdfstring{Reification --- \texttt{only\_enforce\_if}}{Reification --- only\_enforce\_if}}\label{sec:7.3}

Apply a constraint only \emph{when} a Boolean (or a conjunction of
Booleans) is true.

\begin{lstlisting}[language=Python]
load = model.new_int_var(0, 1000, "load")

# Capacity depends on which truck we use.
model.add(load <= 50).only_enforce_if(truck_a)
model.add(load <= 80).only_enforce_if(truck_b)

# Conjunction of literals: enforced only when b1 is true and b2 is false.
model.add(x + y == 10).only_enforce_if([b1, ~b2])
\end{lstlisting}

\lstinline!only_enforce_if! is the
conditional-constraint operator. The constraint is in force when the
literal (or all of the literals in the list) is true; when any of them
is false, the constraint is treated as inactive --- neither violated nor
satisfied, simply absent for that solution. It is the typical way to
express that a constraint should depend on a Boolean.

Not every constraint supports
\lstinline!only_enforce_if!, and the cost of using it
varies considerably across those that do --- cheap on linear and Boolean
constraints, much heavier on the scheduling and routing globals, where a
native mechanism (optional intervals, self-loop literals) is usually the
better expression of conditionality. The OR-Tools documentation lists
which constraints accept it.

\section{Full Reification --- A Boolean Equivalent to a Relation}\label{sec:7.4}

\lstinline!only_enforce_if! from \Cref{sec:7.3} is \emph{half}
reification: it forces a constraint when a literal is true but says
nothing about the false case. Posting it in both directions --- once for
the literal, once for its negation --- gives the full equivalence and
turns a numeric relation (\emph{exceeds threshold}, \emph{is over
budget}) into a Boolean the rest of the model can reason about.

\begin{lstlisting}[language=Python]
over = model.new_bool_var("over_limit")

# over == 1   iff   x >= threshold
model.add(x >= threshold).only_enforce_if(over)
model.add(x <= threshold - 1).only_enforce_if(~over)
\end{lstlisting}

The two \lstinline!only_enforce_if! lines together
establish the equivalence
\(\text{over} \Leftrightarrow (x \ge \text{threshold})\). Once reified,
\lstinline!over! is a first-class Boolean: it can appear
in the objective (\emph{penalty for going over}), in further logical
constraints (\emph{if over the limit, route through warehouse two}), or
anywhere else a literal is accepted.

Any time the model needs to \emph{talk about} whether a numeric relation
holds --- to penalize it, branch on it, combine it with other logic ---
full reification produces the Boolean to talk about.

Often, however, half-reification is enough. If the Boolean's only
purpose is to drive a penalty in the objective --- pay when
\lstinline!over! goes high, no rule when it does not ---
the second \lstinline!only_enforce_if! is unnecessary:
the minimization keeps \lstinline!over! at zero whenever
the relation does not force it to one. Full reification is required only
when the \emph{false} direction itself drives logic --- \emph{if not
over, then route through warehouse two}. Posting both directions when
one would do costs extra constraints and propagation, so drop the half
that is not needed.

\section{\texorpdfstring{\texttt{add\_all\_different}}{add\_all\_different}}\label{sec:7.5}

A list of variables must take pairwise distinct values.

\begin{lstlisting}[language=Python]
model.add_all_different([x1, x2, x3, x4])

# Affine expressions also work — the classic n-queens diagonal trick:
model.add_all_different([queens[i] + i for i in range(n)])
model.add_all_different([queens[i] - i for i in range(n)])
\end{lstlisting}

Mathematically equivalent to \(\binom{n}{2}\) pairwise
\passthrough{\lstinline"!="} constraints, but with a \emph{strictly
stronger} propagator in the domain-reasoning sense: the global form
reasons about the entire list at once --- bipartite matching, in essence
--- and prunes domain values that no pairwise
\passthrough{\lstinline"!="} could discover.

One small note: CP-SAT will auto-lift a cluster of pairwise
\passthrough{\lstinline"!="} constraints into an all-different
propagator when no \lstinline!add_all_different! is
already declared on the same variables, so writing the pairwise form
does not necessarily forfeit global propagation. Declaring both forms on
the same variables disables that auto-lift, which is worth being aware
of.

The constraint accepts not only bare variables but also affine
expressions. The n-queens diagonal trick is the elegant illustration:
queen \(i\) sits in column \lstinline!queens[i]!, on
diagonal \lstinline!queens[i] + i!, and on anti-diagonal
\lstinline!queens[i] - i!. No two queens may share any of
the three, and three \lstinline!add_all_different! calls
express the entire problem.

Typical uses are conflict and uniqueness rules in real models. Time-slot
conflicts where no two tasks may share a start time. Strict orderings,
where the variables encode positions in a permutation and every position
must be distinct. Unique resource assignment --- each task takes its own
machine, each request gets its own identifier, each vehicle a distinct
depot. Frequency assignment where neighboring transmitters must not
share a frequency. Anything that reads \emph{``no two X share a Y''}
fits naturally; sudoku and n-queens are the textbook illustrations of
the same pattern, but the constraint earns its place by appearing in
scheduling, allocation, and routing models far more often than in
puzzles.

\section{\texorpdfstring{\texttt{add\_circuit} --- TSP and Routing}{add\_circuit --- TSP and Routing}}\label{sec:7.6}

A list of \lstinline!(tail, head, literal)! triples; the
chosen arcs (literal = true) must form a single Hamiltonian circuit.

\begin{lstlisting}[language=Python]
arcs = [
    (u, v, model.new_bool_var(f"e_{u}_{v}"))
    for (u, v) in directed_edges
]
model.add_circuit(arcs)
model.minimize(sum(cost[u, v] * lit for (u, v, lit) in arcs))
\end{lstlisting}

The circuit propagator handles subtour elimination internally ---
generating the cuts that exclude shorter cycles on the fly --- so the
model does not need position variables, element constraints, or any of
the manual scaffolding the natural formulation would require.

Several adaptations fall out of the same shape. Optional vertices: add a
self-loop \lstinline!(v, v, skip_v)!; if
\lstinline!skip_v = 1!, the vertex is bypassed. Shortest
path between two endpoints: zero-cost self-loops on every non-endpoint
plus a forced closing arc. Prize-collecting TSP: penalize the skip-loop
literals in the objective, letting the solver decide which vertices are
worth visiting. Multiple tours sharing a depot are covered by
\lstinline!add_multiple_circuit!, the vehicle-routing
variant.

CP-SAT's \lstinline!add_circuit! handles surprisingly
large routing instances --- a few hundred nodes, in many cases, and with
rich side constraints. For \emph{very} large pure TSP --- tens of
thousands of cities --- specialist solvers like Concorde and LKH still
win, as \Cref{sec:5.2} noted.

\section{\texorpdfstring{\texttt{add\_no\_overlap} --- Single-Machine Scheduling}{add\_no\_overlap --- Single-Machine Scheduling}}\label{sec:7.7}

A list of intervals that may not overlap on the time axis.

\begin{lstlisting}[language=Python]
intervals = [
    model.new_fixed_size_interval_var(start[i], duration[i], f"task_{i}")
    for i in range(n_tasks)
]
model.add_no_overlap(intervals)
\end{lstlisting}

Captures a single resource --- machine, room, employee --- that does at
most one thing at a time. Optional intervals (\Cref{sec:6.4}) are respected
automatically: when their presence Boolean is false, they are ignored.

For job-shop scheduling --- multiple machines, each running its own
subset of tasks --- the natural shape is one
\lstinline!add_no_overlap! call per machine:

\begin{lstlisting}[language=Python]
for m in machines:
    model.add_no_overlap([intervals[t, m] for t in tasks if can_run(t, m)])
\end{lstlisting}

The propagator behind \lstinline!add_no_overlap! draws
on \emph{edge finding}, \emph{detectable precedences}, and
\emph{energetic reasoning}. Hand-rolling the same constraint with
pairwise disjunctions on integer start times is correct but tends to be
slower on non-trivial sizes.

\section{\texorpdfstring{\texttt{add\_no\_overlap\_2d} --- 2D Packing}{add\_no\_overlap\_2d --- 2D Packing}}\label{sec:7.8}

The two-dimensional analogue: axis-aligned rectangles in the plane that
may not overlap.

\begin{lstlisting}[language=Python]
x_intervals = [
    model.new_fixed_size_interval_var(xs[i], w[i], f"x_{i}")
    for i in range(n)
]
y_intervals = [
    model.new_fixed_size_interval_var(ys[i], h[i], f"y_{i}")
    for i in range(n)
]
model.add_no_overlap_2d(x_intervals, y_intervals)
\end{lstlisting}

Each rectangle is described by two intervals --- one for its x-extent,
one for its y-extent. The constraint enforces that for every pair of
rectangles, \emph{either} their x-intervals are disjoint \emph{or} their
y-intervals are disjoint, which is exactly the geometric definition of
non-overlap.

Typical uses are 2D bin packing, pallet loading, floor planning, chip
placement, and cutting-stock with rectangular pieces. Optional
rectangles work the same way --- items that may or may not be packed
have their presence Boolean carried by both interval components.

\section{\texorpdfstring{\texttt{add\_cumulative} --- Resource-Capacity Scheduling}{add\_cumulative --- Resource-Capacity Scheduling}}\label{sec:7.9}

The capacitated generalization of
\lstinline!add_no_overlap!. Each interval has a
\emph{demand}; at every moment, the sum of demands of intervals running
at that moment must not exceed a \emph{capacity}.

\begin{lstlisting}[language=Python]
model.add_cumulative(
    intervals=[i0, i1, i2, i3],
    demands=[2, 3, 1, 4],     # paired by index
    capacity=5,
)
\end{lstlisting}

Typical uses are resource-constrained project scheduling (\emph{RCPSP}),
where tasks share a pool of workers, machines, or electricity;
parallel-machine settings with shared capacity, like meeting rooms
holding \(N\) simultaneous sessions; and inventory in a warehouse, where
the ``demand'' of an interval is the units it consumes and the capacity
is the warehouse size.

Capacity itself can be a CP-SAT variable rather than a constant:
declaring \lstinline!capacity! as an integer variable and
minimizing it returns the smallest pool that still completes the work
--- a useful answer to capacity-planning questions.

When the capacity is exactly 1, the specialized
\lstinline!add_no_overlap! exists for the
single-resource case and tends to propagate more strongly than
\lstinline!add_cumulative! on that specialization.

\section{Equality Helpers --- Max, Min, Abs, Multiplication, Division, Modulo}\label{sec:7.10}

The non-linear corners of integer arithmetic ---
\lstinline!max!, \lstinline!min!, \(|x|\),
\(x \cdot y\), \(x \div y\), \(x \mod m\) --- are not legal inside
\lstinline!model.add! directly. CP-SAT exposes them as
auxiliary-variable equalities:

\begin{lstlisting}[language=Python]
m = model.new_int_var(0, 100, "max")
model.add_max_equality(m, [x, y, z])

mn = model.new_int_var(0, 100, "min")
model.add_min_equality(mn, [x, y, z])

abs_x = model.new_int_var(0, 100, "abs_x")
model.add_abs_equality(target=abs_x, expr=x)

xy = model.new_int_var(-1_000_000, 1_000_000, "xy")
model.add_multiplication_equality(xy, [x, y])

q = model.new_int_var(0, 100, "q")
model.add_division_equality(q, x, z)  # integer division: 5 // 2 == 2

r = model.new_int_var(0, 2, "r")
model.add_modulo_equality(r, x, 3)
\end{lstlisting}

The pattern is uniform: declare an auxiliary variable, then call the
matching \lstinline!add_*_equality! to bind it to the
non-linear expression. The auxiliary then participates in further
constraints just like any other integer variable.

Some of these can be sidestepped with linear inequalities ---
\emph{minimizing the max} of values, for instance, needs only
\(m \ge \text{value}_i\) for each \(i\) and
\lstinline!model.minimize(m)!. Reach for the equality
helpers when the relation is genuinely nonlinear in the model.

\section{A Few More, Briefly}\label{sec:7.11}

A handful of constraints we will not dwell on but that are worth
recognizing when the problem calls for them.

\textbf{\lstinline!add_element!} --- binds a target to an
array entry whose index is a decision variable:
\lstinline!target == array[index]!. Useful for lookup
tables, piecewise functions on a discrete domain, and any ``pick the
value at this index'' rule. Array entries may be constants or variables.

\textbf{\lstinline!add_allowed_assignments!} /
\textbf{\lstinline!add_forbidden_assignments!} --- a
tuple of variables must take its joint value from (or avoid) an explicit
list of allowed combinations. Useful for allowed shift patterns,
compatible feature bundles, and truth-table-style rules. Cost grows with
the number of rows; when the rule has sequential structure,
\lstinline!add_automaton! below is often more compact.

\textbf{\lstinline!add_multiple_circuit!} ---
vehicle-routing-style problems where several tours share a depot. Same
triple-list shape as \lstinline!add_circuit!; the
propagator allows multiple disjoint circuits provided they all pass
through the designated depot vertex.

\textbf{\lstinline!add_reservoir_constraint!} --- a
running level (inventory, staff on duty, water in a tank) bounded by
\([\min, \max]\) as discrete \lstinline!+/-! events fire
over time. Inputs are a list of times and a list of per-event level
changes; the constraint enforces that the running sum stays within the
bounds.

\textbf{\lstinline!add_automaton!} --- a sequence of
integer variables must follow a finite-state machine. Useful for
forbidden patterns in a roster (\emph{no more than three night shifts in
a row}), protocol compliance, batch-scheduling state transitions. The
right tool when the rule is naturally stated as ``the sequence is
accepted by this DFA.''

\textbf{\lstinline!add_inverse(v, w)!} --- two
equal-length variable arrays that must be each other's inverses:
\lstinline!v[i] = j! if and only if
\lstinline!w[j] = i!. Used in stable-matching shapes and
any time a problem has two natural representations of the same
combinatorial object that must agree (the modeling pattern is called
\emph{channelling between dual encodings}).

The CP-SAT documentation is the authoritative reference for the corners
of each.
